% SPDX-License-Identifier: MIT
% Copyright (c) 2026 Evgeny Baulin
% Optimization Methods, Seminar 08: one-page cheat sheet (Russian).
% Generated by checks/08. Constrained Optimality and Discrete Modelling.ipynb from web/08/js/content-cheatsheet.js and content/students.json. Do not edit by hand.
% Compile with: latexmk -xelatex -interaction=nonstopmode -halt-on-error Cheatsheet_ru.tex
\RequirePackage{fix-cm}
\documentclass[10pt,a4paper]{article}

\usepackage{amsmath,amssymb}
\usepackage{fontspec}
\setmainfont{cmunrm.otf}[BoldFont=cmunbx.otf, ItalicFont=cmunti.otf, BoldItalicFont=cmunbi.otf]
\setsansfont{cmunss.otf}[BoldFont=cmunsx.otf, ItalicFont=cmunsi.otf, BoldItalicFont=cmunso.otf]
\setmonofont{cmuntt.otf}[BoldFont=cmuntb.otf, ItalicFont=cmunit.otf, BoldItalicFont=cmuntx.otf]
\usepackage{polyglossia}
\setdefaultlanguage{russian}
\setotherlanguage{english}
\newfontfamily\cyrillicfont{cmunrm.otf}[Ligatures=TeX, BoldFont=cmunbx.otf, ItalicFont=cmunti.otf, BoldItalicFont=cmunbi.otf]
\newfontfamily\cyrillicfontsf{cmunss.otf}[Ligatures=TeX, BoldFont=cmunsx.otf, ItalicFont=cmunsi.otf, BoldItalicFont=cmunso.otf]
\newfontfamily\cyrillicfonttt{cmuntt.otf}[BoldFont=cmuntb.otf, ItalicFont=cmunit.otf, BoldItalicFont=cmuntx.otf]
\PolyglossiaSetup{russian}{hyphenmins={2,3}}

\usepackage[a4paper,left=8mm,right=8mm,top=7mm,bottom=11mm,footskip=5.5mm]{geometry}
\usepackage{xcolor}
\usepackage{multicol}
\usepackage{enumitem}
\usepackage{ragged2e}
\usepackage{fancyhdr}
\usepackage[colorlinks=true,linkcolor=hseblue,urlcolor=hseblue]{hyperref}

\definecolor{hseblue}{RGB}{16,45,105}
\definecolor{hseaccent}{RGB}{198,40,40}
\definecolor{hsegray}{RGB}{108,114,126}
\definecolor{hselight}{RGB}{240,243,248}

\hypersetup{pdftitle={Шпаргалка к семинару 8: Условная оптимальность и дискретное моделирование}, pdfauthor={Евгений Баулин}}

\pagestyle{fancy}
\fancyhf{}
\renewcommand{\headrulewidth}{0pt}
\fancyfoot[L]{\scriptsize\color{hsegray}\textcopyright\ 2026 Евгений Баулин, \href{https://t.me/tarakan_tuc}{t.me/tarakan\_tuc}}
\fancyfoot[R]{\scriptsize\color{hsegray}Лицензия MIT}

\setlength{\parindent}{0pt}
\setlength{\parskip}{0pt}
\setlength{\columnsep}{4.5mm}
\setlength{\columnseprule}{0.3pt}
\renewcommand{\columnseprulecolor}{\color{hselight}}
\setlength{\emergencystretch}{2em}
\setlength{\RaggedRightRightskip}{0pt plus 4em}
\binoppenalty=2000
\relpenalty=500
\setlength{\abovedisplayskip}{2pt}
\setlength{\belowdisplayskip}{2pt}
\setlist[itemize]{leftmargin=7pt, labelsep=3pt, itemsep=1.2pt, topsep=1pt, parsep=0pt, partopsep=0pt,
  label={\color{hseaccent}\scriptsize\textbullet}}

\makeatletter
\newcommand{\cssec}[1]{\par\addvspace{3.5pt}\noindent
  \colorbox{hselight}{\parbox{\dimexpr\linewidth-2\fboxsep}{\RaggedRight\hyphenpenalty=10000\color{hseblue}\bfseries\strut #1}}%
  \par\nobreak\vspace{1pt}\@afterheading}
\makeatother
\clubpenalty=10000
\widowpenalty=10000

\begin{document}
\thispagestyle{fancy}

\noindent
\begin{minipage}[b]{0.6\textwidth}
  \RaggedRight\hyphenpenalty=10000
  {\color{hseaccent}\bfseries\small Шпаргалка к семинару 8}\\[1pt]
  {\color{hseblue}\bfseries\Large Условная оптимальность\\ и дискретное моделирование}
\end{minipage}\hfill
\begin{minipage}[b]{0.39\textwidth}
  \raggedleft\footnotesize
  {\bfseries Методы оптимизации}\\
  НИУ ВШЭ, факультет компьютерных наук, 2026--2027\\
  Лектор курса: Андрей Игнатов\\
  Автор: Евгений Баулин, Telegram \href{https://t.me/tarakan_tuc}{@tarakan\_tuc}
\end{minipage}

\vspace{2pt}
{\color{hseblue}\rule{\textwidth}{0.8pt}}
\vspace{-4pt}

\fontsize{9}{10.6}\selectfont
\begin{multicols}{3}
\RaggedRight


% ---- kkt
\cssec{Условия ККТ}
\begin{itemize}
  \item $L=f+\sum\lambda_ig_i+\sum\nu_jh_j$, $g_i\le0$, $\lambda_i\ge0$, $\nu_j$ любого знака.
  \item Стационарность $\nabla_xL=0$, допустимость, $\lambda\ge0$, $\lambda_ig_i=0$.
  \item Необходимы при условии регулярности (LICQ или Слейтер для выпуклых задач); достаточны для выпуклой задачи.
  \item Строго выполнено $\Rightarrow\lambda_i=0$; из точности не следует $\lambda_i>0$.
\end{itemize}

% ---- cases
\cssec{Перебор активных множеств}
\begin{itemize}
  \item Выберите $A$: $g_i=0$ на $A$, $\lambda_i=0$ вне $A$; решите систему.
  \item Отбросьте, если какое-то $g_i>0$ или $\lambda_i<0$.
  \item Выпуклая задача: уцелевший кандидат глобален. Невыпуклая: сравните всех уцелевших.
\end{itemize}

% ---- dual
\cssec{Двойственность и цены}
\begin{itemize}
  \item $q=\inf_xL$ вогнута; $q(\lambda,\nu)\le f(x)$ для допустимой $x$ и $\lambda\ge0$.
  \item Выпуклость + строго допустимая точка: $d^*=p^*$, и максимум двойственной задачи достигается.
  \item $g_i\le b_i$: $\partial p^*/\partial b_i=-\lambda_i^*$. Для $\ge$ знак меняется.
  \item $f(x)-q(\lambda)\ge0$ показывает, насколько допустимая $x$ может быть далека от оптимума.
\end{itemize}

% ---- milp
\cssec{Целочисленные модели}
\begin{itemize}
  \item Решения $\to$ переменные, правила $\to$ ограничения, цель $\to$ целевая функция.
  \item Релаксация: в максимизации — оценка сверху, в минимизации — снизу; округление — не метод.
  \item $x\le Mz$ включает постоянные затраты; $z_A\le z_B$ — «если A, то B»; $\sum z_i\le k$ ограничивает выбор.
  \item Берите $M$ как можно меньше: те же целые решения, более точная релаксация.
\end{itemize}

\end{multicols}
\end{document}
